\documentclass[12pt,a4paper]{article}
\usepackage[utf8]{inputenc}
\usepackage[T1]{fontenc}
\usepackage[margin=2.5cm]{geometry}
\usepackage{graphicx}
\usepackage[colorlinks=true,
            linkcolor=black,
            urlcolor=blue,
            citecolor=black]{hyperref}
\usepackage{array}
\usepackage{booktabs}
\usepackage{tabularx}
\usepackage{enumitem}
\usepackage{setspace}
\usepackage{titlesec}
\usepackage{parskip}
\usepackage{amsmath}
\usepackage{tcolorbox}
\usepackage{lscape}

\geometry{margin=1.5cm}
\newcolumntype{C}[1]{>{\raggedright\arraybackslash}p{#1}}

\newcommand{\titrule}{\rule[0.5ex]{\linewidth}{1pt}}
\titleformat{\section}{\large\bfseries}{}{0em}{}[\titrule]
\titleformat{\subsection}{\normalsize\bfseries}{}{0em}{}

\setlength{\parindent}{0pt}
\setlength{\parskip}{6pt}
\onehalfspacing

\title{\textbf{Assignment 1: Lean Startup}\\
       \large StudentBridge -- Dual-Tier B2B Managed Project \& R\&D Sandbox\\[0.5em]
       \normalsize 2DV614}
\author{Samuel Fredric Berg}
\date{\today}

\begin{document}

\maketitle
\thispagestyle{empty}
\newpage
\tableofcontents
\newpage

% ============================================================
\section{Business Idea}

\textbf{StudentBridge} is a B2B managed marketplace that connects regional organizations in Småland, Sweden, with highly vetted university students for structured, 10-hour digital, analytical, and software micro-projects. Rather than operating as a passive job board where companies must manually sort through resumes, StudentBridge acts as an active intermediary, scoping coordinator, and contract execution guarantor.

As a platform, StudentBridge functions structurally as a \textbf{two-sided marketplace} with a distinct operational and financial asymmetry. The \textbf{paying consumers} on the demand side of the platform are the regional corporate clients and SMEs who fund the project blocks to clear internal backlogs. Conversely, the university students represent the \textbf{supply side of the ecosystem}, signing up to the platform not as financial buyers, but as an allocatable, highly vetted resource workforce looking to gain vital client-facing portfolio experience before graduation. Consequently, the platform must execute two separate acquisition strategies: a premium B2B sales pipeline targeting corporate capital, and an organic campus placement pipeline focused on professional career acceleration.

Market research indicates that regional economic friction stems primarily from a structural mismatch of time and specialized resource allocation rather than a lack of foundational interest. On one side, established software engineering and IT enterprises face severe developer bandwidth shortages. Data from larger software companies shows that due to internal human resource bottlenecks, engineering teams can only afford to actively pursue roughly 15\%--20\% of incoming product opportunities, customer feature requests, or speculative proof-of-concept (PoC) initiatives. The remaining 80\%--85\% of secondary commercial concepts are backlogged or discarded due to high engineering costs.

Conversely, small and medium-sized enterprises (SMEs) face an absolute resource and departmental constraint. Unlike enterprise peers, regional SMEs frequently operate without dedicated internal IT or digital marketing business units, making small, routine digital updates difficult to execute efficiently without contracting expensive full-service firms. Concurrently, highly capable university graduates frequently exit their academic tracks with strong theoretical competence but face a structural experience gap when competing for entry-level positions.

To accommodate these contrasting operational realities, StudentBridge deploys a dual-tier business offering tailored to explicit corporate needs:

\begin{itemize}
    \item \textbf{The SME Operational Tier:} Built for resource-constrained small businesses. These companies leverage StudentBridge to clear immediate, minor digital overhead tasks (e.g., setting up localized web pages, formatting customer databases, or designing social media collateral). StudentBridge handles the entire pipeline management, removing any need for technical overhead from the SME owner.
    \item \textbf{The Enterprise R\&D Tier:} Built specifically for established software firms needing rapid concept verification. These companies use StudentBridge as an external, low-risk sandbox. Vetted Computer Science and Business students are deployed to build rapid interactive mockups, execute API verification tests, or write standalone documentation to determine if backlogged ideas warrant deeper internal investment.
\end{itemize}

To guarantee delivery standards to paying corporate entities, the platform enforces a strict, multi-stage \textbf{Workforce Vetting and Quality Protocol}. Onboarding as an allocatable resource is governed by baseline academic criteria: students must verify active enrollment and a proven history of decently performing their academic obligations (e.g., meeting credit milestones and passing core prerequisites).

Once accepted into the active registry, the vetting shifts from academic verification to historical platform merit. Subsequent project allocations are dynamically determined by verified professional track records, comprising a portfolio of previously completed micro-projects and direct performance reviews from client companies highlighting execution quality, technical precision, and communication.

The corporate client pays a flat fee per 10-hour project block. StudentBridge holds the funds in escrow, releasing a highly competitive milestone payout and a verified \emph{Certificate of Work} to the student only upon verified delivery against objective ``definition-of-done'' metrics.

To protect the ecosystem's integrity from behavioral liabilities, the platform enforces a strict, legally binding \textbf{Graduated Accountability and Anti-Ghosting Framework}. If an active user (student or corporate client) abruptly halts communication or abandons an active project sprint for more than 48 business hours, a multi-tiered enforcement track is automatically triggered:
\begin{itemize}
    \item \textbf{First Infraction:} The user faces an immediate temporary profile freeze. Students immediately forfeit their escrowed milestone payout, and a flat 1,000 SEK breach-of-contract penalty fee is deducted from past or future platform earnings. Corporate clients are assessed a matching 1,500 SEK liquidation penalty.
    \item \textbf{The Interval Recidivism Ban:} If a second ghosting or communication default infraction occurs within any rolling 12-month interval, it triggers an irreversible system escalation. The offending user's account is permanently deactivated, and they face a lifetime ban from utilizing any StudentBridge corporate services or workforce placement frameworks.
\end{itemize}

% ============================================================
\newpage
\begin{landscape}
\section{Lean Model Canvas}
\thispagestyle{empty}

\tcbset{colback=white, colframe=black, arc=0mm, boxrule=0.5pt, top=4pt, bottom=4pt, left=4pt, right=4pt}
\footnotesize

\setlength{\tabcolsep}{4pt}
\renewcommand{\arraystretch}{1.2}

\begin{tabular}{|C{5.2cm}|C{5.2cm}|C{5.2cm}|C{5.2cm}|C{5.2cm}|}
\hline
\textbf{Problem} & \textbf{Solution} & \textbf{Unique Value Proposition} & \textbf{Unfair Advantage} & \textbf{Customer Segments} \\
\hline
1. University graduates face an industry experience gap upon entering the workforce.\newline
2. Resource-constrained SMEs lack dedicated internal bandwidth to execute minor, necessary digital updates.\newline
3. Enterprise tech firms lack low-risk environments to run rapid proofs-of-concept (PoC) on backlogged ideas. &
A dual-tier B2B platform supplying structured, 10-hour managed sprints backed by a strict anti-ghosting penalty clause and escrow guarantees. &
\textbf{For SMEs:} Hands-off resolution of digital operational bottlenecks.\newline
\textbf{For Tech Firms:} Low-cost R\&D sandboxing to validate backlogged concepts risk-free. &
Institutional integration with Linnaeus University; proprietary dual-tier project templates; contractually enforced user accountability rules. &
\textbf{Two-Sided Marketplace Split:}\newline
\textbf{1. Paying Consumers (Demand):} Local resource-constrained SMEs and large Småland tech firms allocating budget for R\&D validation.\newline
\textbf{2. Experiential Workforce (Supply):} Vetted CS and Business students signing up to gain allocatable, resume-building project experience. \\
\hline
\textbf{Existing Alternatives} & \textbf{Key Metrics} & \textbf{High-Level Concepts} & \textbf{Channels} & \textbf{Early Adopters} \\
\hline
Traditional freelance platforms (Upwork, Fiverr), local IT consultancies, or unmanaged university capstone projects. &
Project completion rate; ghosting/default rate; corporate repeat utilization; student application-to-acceptance ratio. &
"Managed micro-consulting for the overlooked middle market." &
\textbf{B2B:} Tech hub networks (Videum Science Park), Chamber of Commerce.\newline
\textbf{Students:} University career portals, campus posters, union partnerships. &
Non-technical local business owners in Kronoberg; software engineering juniors/seniors looking for local portfolio pieces. \\
\hline
\multicolumn{2}{|C{10.64cm}|}{\textbf{Cost Structure}} & \multicolumn{3}{c|}{\textbf{Revenue Streams}} \\
\hline
\multicolumn{2}{|C{10.64cm}|}{Student escrow payouts (2,000--3,000 SEK/project), enterprise NDA frameworks, dispute resolution overhead, manual auditor re-scoping costs (1,200 SEK per student breach incident), Stripe B2B escrow fees, automated vetting tools.} &
\multicolumn{3}{c|}{\parbox{16.0cm}{\vspace{3pt}
\textbf{SME Tier Fee:} 2,500 SEK per project (500 SEK platform margin).\newline
\textbf{Enterprise Tier Fee:} 4,500 SEK per advanced sprint (1,500 SEK platform margin).\newline
\textbf{Recidivism Penalties:} Strike 1 ghosting fees (1,000 SEK student / 1,500 SEK company); Strike 2 results in immediate structural expulsion.\vspace{3pt}}} \\
\hline
\end{tabular}
\end{landscape}

\newpage

% ============================================================
\section{Assumptions and Hypotheses}

The foundation of the StudentBridge model relies upon addressing systemic resource mismatches through the following prioritized, segmented assumptions:

\begin{enumerate}[label=\textbf{A\arabic*}.]
    \item \textbf{(Symmetrical Two-Sided Liquidity)} StudentBridge can scale efficiently by simultaneously converting regional companies into paying transactional consumers while consistently maintaining an active, un-monetized registry of students willing to sign up as an allocatable resource.
    \item \textbf{(Enterprise Low-Risk Scoping)} Large technology firms possess expansive backlogs of innovative feature requests and concepts, but lack a low-risk, rapid external environment to execute proofs-of-concept (PoC) to determine if these initiatives warrant deeper corporate investment.
    \item \textbf{(SME Resource Constraints)} Local small and medium enterprises (SMEs) face operational digital overhead (e.g., localized web presence, data formatting) but lack the internal technical departments or dedicated resource allocation to justify hiring expensive, full-service IT consultancies.
    \item \textbf{(Acceptance of Accountability Clauses)} Both corporate clients and student applicants will willingly accept contract terms that impose escalating financial penalties and progressive platform expulsion for unresponsiveness or ghosting to maintain system trust.
\end{enumerate}

% ============================================================
\section{Validation of the Three Most Important Assumptions}

\subsection{A1 -- Two-Sided Marketing and Demand Segment Breakdown}

\textbf{Method:} Dual marketing outreach tests were run in Kronoberg County over a two-week validation cycle to verify if both sides of the marketplace could be cost-effectively acquired. On the corporate side, direct email proposals were sent to 20 larger IT/software firms (R\&D sandbox) and 20 small manufacturing/service companies with limited technical departments (managed digital fixes). On the student workforce side, digital flyer campaigns were deployed in Linnaeus University computer science and business group channels. \textit{(For the exact raw text of these outreach templates and experimental design specifications, see Appendix~\ref{app:outreach}).}

\textbf{Results:}
\begin{itemize}
    \item \textbf{Corporate Demand Side:} 14 larger tech firms responded; 11 confirmed that the 15\%--20\% resource bottleneck is highly accurate, and 8 stated they would pay 4,500 SEK to pass validation tasks to students. For smaller SMEs, 9 responded; 6 stated they would pay 2,500 SEK for a standalone web/data update, provided StudentBridge manages the student entirely.
    \item \textbf{Student Workforce Supply Side:} The campaign attracted 95 qualified student sign-ups within 7 days. This confirmed that the supply side can be rapidly aggregated through low-cost organic university career channels. \textit{(The daily registration growth metrics are charted in Appendix~\ref{app:metrics}).}
\end{itemize}

\textbf{Conclusion:} The assumption is \emph{strongly supported}. The market demands separate corporate offerings, and student interest is highly responsive to localized campus marketing.

\subsection{A2 -- Acceptance of Strict Accountability and Ghosting Clauses}

\textbf{Method:} The 95 surveyed students and the 14 engaged corporate product managers were presented with our mandatory contract framework containing the multi-tier enforcement framework: ``A first communication breach during a sprint triggers a 1,000 SEK student fee (1,500 SEK corporate fee), and a second infraction within a 12-month interval results in an immediate permanent lifetime ban.'' \textit{(The precise survey questionnaire architecture is documented in Appendix~\ref{app:surveys}).}

\textbf{Results:}
\begin{itemize}
    \item \textbf{Student Feedback:} 89 out of 95 students (93.6\%) endorsed the tightened rules, noting that a premium 200--300 SEK/hour equivalent rate justifies strict professional guidelines. They reported that the rolling 12-month ban threshold heavy protects their own workforce reputation by filtering out unreliable peers.
    \item \textbf{Corporate Feedback:} 14 out of 14 product managers (100\%) aggressively supported the graduated ban framework. They noted that student unreliability and ghosting was their primary fear when considering working with unmanaged freelancers, and that a permanent expulsion mechanism heavily increased their willingness to integrate StudentBridge into their core engineering workflows.
\end{itemize}

\textbf{Conclusion:} The assumption is \emph{strongly supported}. The escalating accountability clause acts as a powerful B2B trust-building asset rather than a point of transaction friction.

\subsection{A3 -- Multi-Tier Technical Scoping Feasibility}

\textbf{Method:} We ran detailed diagnostic scoping sessions with 3 small business owners (SMEs) and 3 product managers from large software firms to evaluate their respective operational backlogs.

\textbf{Results:} Standardized agile templates successfully split tasks into predictable 10-hour segments across both distinct groups:
\begin{itemize}
    \item \textbf{SME Tier Tasks:} Isolated into straightforward operational goals, such as building a responsive landing page using WordPress or cleaning up historical client email spreadsheets.
    \item \textbf{Enterprise Tier Tasks:} Isolated into conceptual proof-of-concepts, such as creating an isolated React module to test interactive UI animations, or writing a Python script to verify response data from an external corporate API.
\end{itemize}

\textbf{Conclusion:} The assumption is \emph{supported}. Diverse corporate needs can be extracted into highly organized 10-hour project blocks by using tier-specific scoping guidelines.

% ============================================================
\section{Findings Summary}

\subsection{Key Marketplace Takeaways}
The validation loop confirmed that constructing StudentBridge as a segmented, highly contractually protected two-sided marketplace optimizes regional market viability. Key findings include:
\begin{itemize}
    \item \textbf{Dual Marketing Mandate:} Marketing cannot be generalized. Corporate marketing must focus on risk reduction, ROI, and opportunity reclamation. Student marketing must focus on high-trust competitive payouts, resume expansion, and institutional career credibility.
    \item \textbf{Accountability and Quality Control:} Trust is the fundamental product of StudentBridge. Implementing financial penalties and expulsion mechanics for ghosting mitigates marketplace degradation, matching the institutional standards required by enterprise tech firms.
    \item \textbf{Divergent Sales Cycles:} Selling to the Enterprise tier involves navigating institutional non-disclosure and intellectual property (IP) legal clearances. The SME tier sales cycle is purely transactional and focused on immediate problem-solving.
\end{itemize}

% ============================================================
\section{Scenario Analysis}

\subsection{Cost Structure and Financial Metrics}

The baseline financial model utilizes a diversified transaction matrix consisting of 70\% Standard SME Projects (500 SEK platform margin) and 30\% Advanced Enterprise Sprints (1,500 SEK platform margin). This calculation assumes zero net revenue from penalty fees, treating them purely as behavioral deterrents that offset platform damage control overhead.

\textbf{Weighted Average Net Revenue per Transaction:}
\[
  \text{Blended Platform Margin} = (0.70 \times 500\text{ SEK}) + (0.30 \times 1,500\text{ SEK}) = 800\text{ SEK}
\]

\textbf{Monthly Fixed Costs (MVP Phase):}

\begin{tabularx}{\textwidth}{Xr}
\toprule
Item & Cost (SEK/month) \\
\midrule
Platform Infrastructure (Hosting, Escrow Gateways, Vetting Tools) & 1,200 \\
B2B Corporate Marketing, Tech Park Network Dues, Sales Assets & 2,000 \\
University Student Recruitment (Campus Marketing, Printing, Portals) & 1,000 \\
Part-time Code Quality Auditor / Student Assessor (Stipend) & 8,000 \\
Employer Taxes \& Social Contributions (\emph{Arbetsgivaravgift} \textasciitilde 31.42\%) & 2,514 \\
Corporate Legal Fees (Pre-cleared Enterprise NDA \& Anti-Ghosting Terms) & 1,500 \\
\midrule
\textbf{Total Monthly Fixed Costs} & \textbf{16\,214} \\
\bottomrule
\end{tabularx}

\textbf{Break-even Calculation:}
\[
  \text{Break-even Volume} = \frac{16\,214 \text{ SEK}}{800 \text{ SEK}} \approx 21 \text{ transactions/month}
\]
Due to the premium margins collected from the Enterprise tier, the business requires only 21 total completed assignments per month across the regional ecosystem to offset operational overhead.

\textbf{Total 500\,000~SEK Runway (Zero Growth Baseline):}
\[
  \text{Runway} = \frac{500\,000}{16\,214} \approx 30.8 \text{ months}
\]

\subsection{Scenario 1 -- Base Case (Believed Symmetrical Growth)}

The platform balances its two-sided marketplace demands by deploying targeted campus outreach alongside integrations within regional hubs like Videum Science Park. Growth remains balanced between student labor and corporate project listings.

\begin{tabularx}{\textwidth}{Xrrrr}
\toprule
Month & Total Projects & SME Tier (70\%) & Enterprise Tier (30\%) & Net Cash Flow (SEK) \\
\midrule
1  &   8 &   6 &   2 & -9\,814 \\
3  &  18 &  13 &   5 & -1\,814 \\
6  &  35 &  25 &  10 & +11\,786 \\
12 &  75 &  53 &  22 & +43\,786 \\
\bottomrule
\end{tabularx}

Net profitability is attained during Month 4 or 5. Total runway consumed from the 500,000 SEK grant prior to achieving regular self-sustainability is under 28,000 SEK, leaving the remaining funding available for tech scaling.

\subsection{Scenario 2 -- Pessimistic (Asymmetrical Marketplace Friction)}

The platform suffers from a classic two-sided marketplace failure: student workforce recruitment is highly successful, but corporate B2B sales cycles drag due to corporate legal hesitation regarding anti-ghosting terms. The platform averages a flat volume of only 10 total projects per month through year one, leaving a surplus of underutilized students on the platform.

Monthly platform revenue is $10 \times 800 = 8,000$ SEK against fixed outlays of 16,214 SEK, producing a monthly loss of 8,214 SEK. Over a one-year timeframe, total capital burn reaches 98,568 SEK. Protected by the 500,000 SEK incubator foundation, the startup can comfortably sustain operations in this slow state for more than 4.5 years, giving the founders extensive time to iterate their B2B sales methodology.

\textbf{Mitigation Strategy:} Shift corporate outreach to focus heavily on pre-vetted institutional partners or chamber members who can bypass standard legal review bottlenecks.

\subsection{Scenario 3 -- Optimistic (Viral Two-Sided Network Effects)}

Spurred by compounding engineering talent shortages, local software enterprises swiftly offload backlogs. Simultaneously, word-of-mouth on campus creates a viral pool of top-tier student applicants, completely eliminating acquisition costs for the student workforce side. Enterprise-level projects grow to represent 50\% of total platform volume by Month 3.

By Month 3, total volume hits 40 projects (20 SME / 20 Enterprise), yielding a combined net income of $(20 \times 500) + (20 \times 1,500) = 40,000$ SEK, resulting in an immediate net profit of +23,786 SEK/month. By Month 12, the system handles 150 projects monthly, creating over 150,000 SEK in net monthly profit. The primary 500,000 SEK grant remains completely untouched as an equity reserve.

\subsection{Hiring and Infrastructure Scenarios (Task 6d \& 6g)}

\begin{itemize}
    \item \textbf{Personnel Scaling (6d):} The MVP lifecycle operates without dedicated UX or Copywriters; the original student founders split operational management---one founder handles student community acquisition and vetting, while the other drives corporate business development. When metrics scale beyond 50 monthly deliveries (Scenario 1, Month 9), a full-time Senior Code Auditor will be hired to handle enterprise code reviews, followed by an automated customer support manager to oversee dispute resolutions.
    \item \textbf{Operational Consequences of Scale (6g):} Scaling the marketplace multiplies the operational friction of user disputes and communication defaults. When a student default occurs, it introduces a direct operational recovery cost to StudentBridge: the platform must absorb approximately 1,200 SEK in labor overhead per incident to have our internal Code Quality Auditor manually review the abandoned repository, run damage control with the enterprise corporate contact, and instantly re-allocate a backup vetted student to finish the 10-hour sprint. As volume scales, an automated logging and performance monitoring framework must be deployed to record communication timestamps, project delivery benchmarks, and historical company reviews. This automated data tracking serves as the objective source-of-truth used to algorithmically compute individual student matching priority, compile peer-review portfolios for future project allocation, and systematically issue graduated interval ghosting penalties without requiring manual intervention from the core startup team.
\end{itemize}

% ============================================================
\section{Conclusion}

Evolving StudentBridge into a contractually secure, two-sided B2B platform directly aligns with the varied economic realities of Småland. By identifying students as an active, professional workforce supply line and corporations as a specialized demand pipeline, the business avoids common marketplace liquidity risks. The deployment of a strict anti-ghosting and expulsion framework structurally de-risks execution, providing resource-constrained SMEs with secure operational relief and larger software corporations with a robust R\&D validation pipeline. This dual-market focus combined with clear behavioral safeguards reduces break-even requirements to a highly attainable 21 monthly transactions. Backed by a 500,000 SEK grant, StudentBridge possesses the financial resilience and structural clarity required to establish itself as a trusted regional corporate partner. The immediate next action step is to finalize the distinct legal framework templates for both the student and enterprise terms of use.

\newpage
\appendix
\titleformat{\subsection}{\normalsize\bfseries}{\thesubsection\quad}{0em}{}
\renewcommand{\thesubsection}{\Alph{section}.\arabic{subsection}}

\setcounter{page}{1}
\renewcommand{\thepage}{\Alph{page}}
% ============================================================

\section{Appendix: Raw Validation Evidence}

\subsection{B2B Corporate Outreach Email Scripts}
\label{app:outreach}

Below are the anonymized cold outreach templates deployed during the 14-day validation phase to test corporate pricing elasticity and structural appetite.

\begin{tcolorbox}[title=Template A: Enterprise R\&D Sandbox Testing (Sent to 20 IT Firms), colback=white, colframe=black, arc=0mm]
\small
\textbf{Subject:} Structural Backlog Validation Strategy via LNU Tech Sandbox\\
\\
Dear Product Lead, \\
\\
Data indicates that regional engineering resource limitations force mid-to-large software enterprises in Småland to backlog roughly 80\% of incoming feature requests and speculative PoC concepts. \\
\\
StudentBridge is validating an isolated, B2B managed sandbox. We deploy vetted, top-performing Computer Science seniors from Linnaeus University to execute rapid 10-hour technical sprints (e.g., standalone API testing, interactive UI component modeling, documentation) under a strict contractual anti-ghosting penalty. \\
\\
We are testing a flat corporate tier pricing model of \textbf{4,500 SEK per 10-hour sprint block}, managed entirely through third-party milestone escrow. Would your product team be willing to utilize this framework to validate an un-resourced backlog item? \\
\\
Best regards,\\
The StudentBridge Validation Team
\end{tcolorbox}

\begin{tcolorbox}[title=Template B: SME Digital Operations Testing (Sent to 20 Non-Technical Businesses), colback=white, colframe=black, arc=0mm]
\small
\textbf{Subject:} Hands-Off Resolution of Digital Overhead Tasks\\
\\
Dear General Manager, \\
\\
Local small businesses frequently experience operational delay due to routine digital, database, or media maintenance, while lacking the internal technical bandwidth to hire full-time web consultancies. \\
\\
StudentBridge provides a fully managed, hands-off execution pipeline. We coordinate pre-screened university business and technical students to resolve minor digital bottlenecks (e.g., localized WordPress updates, data entry automation, inventory script setups) in structured 10-hour blocks.\\
\\
The business pays a flat fee of \textbf{2,500 SEK per project block}. StudentBridge holds the funds in escrow and completely handles quality control and student management. If this managed service would resolve immediate overhead, please respond to schedule an interview.\\
\\
Best regards,\\
The StudentBridge Validation Team
\end{tcolorbox}

\newpage
\subsection{Workforce Survey Questionnaires \& Contract Framework}
\label{app:surveys}

The following questions were served to the cohort of 95 students and 14 product managers via digital forms to validate assumption \textbf{A2} (Accountability Clauses).

\begin{enumerate}
    \item \emph{For Students:} ``To ensure corporate buyers treat platform work seriously, StudentBridge enforces a 48-hour communication default clause. If you accept a sprint block and cease communication for over 48 business hours without prior approval, you forfeit your escrow milestone and face a 1,000 SEK contract penalty fee. A second default results in a lifetime platform deactivation. Do you consider this threshold fair given a premium 200--300 SEK/hour equivalent compensation rate?''
    \item \emph{For Corporate Product Managers:} ``Our primary deterrent against student unreliability is an escalating 12-month interval recidivism ban. If a student defaults on communication or deadlines twice within a rolling year, they are permanently expelled from our corporate delivery pipelines. Does this explicit structural mechanism eliminate your risk hesitation regarding student talent onboarding?''
\end{enumerate}

\subsection{Student Workforce Acquisition Metrics}
\label{app:metrics}

Organic student supply velocity was captured over a 7-day tracking window following a single flyer drop across local institutional communication servers.

\begin{table}[ht]
\centering
\small
\caption{Daily Cumulative Student Onboarding Metrics (7-Day Window)}
\begin{tabularx}{0.8\textwidth}{Xrr}
\toprule
Experimental Day & New Registry Accounts & Verified Credit Performance \\
\midrule
Day 1 & 18 & 14 \\
Day 2 & 42 & 38 \\
Day 3 & 61 & 55 \\
Day 4 & 73 & 68 \\
Day 5 & 82 & 76 \\
Day 6 & 89 & 83 \\
Day 7 & \textbf{95} & \textbf{89} \\
\bottomrule
\end{tabularx}
\end{table}

This rapid progression indicates that the supply curve is highly elastic and responsive to professional career-building messaging, minimizing corporate customer-to-resource alignment delays during growth phases.

\end{document}